\section{Methodology: SVD-Lin Attention}
\subsection{Core Idea}
Given the hidden state $H \in \mathbb{R}^{N \times d}$ at a particular layer, we first compute its truncated SVD to obtain the top-$r$ left singular vectors $U_r$. This matrix provides an orthonormal basis that captures the principal directions of the sequence. We then use $U_r$ to project the keys and values:
\[
\tilde{K} = U_r^\top K \in \mathbb{R}^{r \times d}, \quad
\tilde{V} = U_r^\top V \in \mathbb{R}^{r \times d}.
\]
The attention scores are computed as:
\[
S = Q \tilde{K}^\top \in \mathbb{R}^{N \times r},
\]
followed by softmax row-wise:
\[
A = \text{softmax}(S / \sqrt{d_k}),
\]
and finally the output is:
\[
O = A \tilde{V} \in \mathbb{R}^{N \times d}.
\]

\subsection{Why Project Keys and Values, Not Queries?}
Projecting queries would reduce the output dimension, losing information. By projecting keys and values, we compress the memory (the set of tokens to attend to) while keeping the queries intact. This is analogous to reducing the number of context vectors in a memory network.

\subsection{Mathematical Derivation of Complexity}
Let $H$ be of size $N \times d$. Computing the randomized SVD costs $O(N d \log r + (N+d) r^2)$. Then:
\begin{itemize}
    \item Project $K$: $U_r^\top K$ costs $O(N r d)$ (matrix multiplication).
    \item Project $V$: same $O(N r d)$.
    \item Compute $Q \tilde{K}^\top$: $O(N r d)$.
    \item Softmax and output: $O(N r d)$.
\end{itemize}
Total per layer: $O(N d \log r + (N+d) r^2 + N r d)$. For typical values $d=768$, $r=128$, $N \gg d$, the dominant term is $O(N r d)$, which is linear in $N$.

\subsection{Handling Multi-Head Attention}
In multi-head attention, we have $h$ heads, each with dimension $d_h = d / h$. The hidden state $H$ is shared across heads, so we compute $U_r$ once from $H$ and use the same projection for all heads. This avoids per-head SVD overhead. The Q, K, V projections are head-specific, but $U_r$ is shared. The final output is concatenated and linearly projected as in standard BERT.

\subsection{Fallback to Full Attention}
If $N \le r$, there is no compression benefit, and we simply use the standard attention mechanism. In practice, for GLUE tasks with maximum length 128, we set $r \ge 64$ so that fallback rarely occurs (only for very short sequences).

\subsection{Algorithm \& Pseudo-Code}
\begin{lstlisting}[language=Python, caption={SVD-Lin Attention Forward Pass}]
# SVD-Lin Attention Forward Pass
# Input: H in R^{B x N x d}, rank r, optional mask M
Q = H @ W_Q; K = H @ W_K; V = H @ W_V
Q, K, V = reshape_heads(Q, K, V)  # [B, h, N, d_h]

if N > r:
    # 1. Compute left singular vectors U_r
    U, Sigma, Vt = torch.svd_lowrank(H, q=r)
    U_r = U[:, :, :r].unsqueeze(1)  # [B, 1, N, r]

    # 2. Compress Keys and Values via U_r projection
    K_tilde = torch.einsum('bhnk,bhkd->bhrd', U_r, K)
    V_tilde = torch.einsum('bhnk,bhkd->bhrd', U_r, V)

    # 3. Compute Low-Rank Attention Scores
    S = torch.einsum('bhmd,bhrd->bhmr', Q, K_tilde) / sqrt(d_h)
    if M is not None:
        S = S + M.unsqueeze(-1)
    A = torch.softmax(S, dim=-1)
    O = torch.einsum('bhmr,bhrd->bhmd', A, V_tilde)
else:
    # Standard O(N^2) full attention fallback
    O = standard_scaled_dot_product(Q, K, V, M)

return output_projection(concat_heads(O))
\end{lstlisting}

\subsection{Integration into BERT}
We replace each `BertSelfAttention` layer in the HuggingFace implementation with our `SVDLinearAttention`. The rest of the BERT architecture (layer norm, feed-forward networks, residuals) remains unchanged. We initialize the linear projections $W_Q, W_K, W_V$ from the pre-trained BERT weights to preserve the original model's knowledge. The SVD projection does not introduce new trainable parameters, so the model size is identical to BERT.

\subsection{Computational Considerations on M1 Mac}
Apple's M1 chip includes a 16-core Neural Engine and a powerful GPU accessible via Metal Performance Shaders (MPS). PyTorch's MPS backend supports most tensor operations, including matrix multiplication and SVD. However, `torch.svd\_lowrank` is not yet natively implemented for MPS; we fall back to CPU for the SVD computation and then move tensors back to MPS. This overhead is acceptable because SVD is only called once per layer per batch element, and the cost is small relative to the attention computation. We also use mixed precision (float16) where possible to accelerate matrix multiplications.